\section{Прикладные особенности}
\label{sec:AppliedAspects}
\tikzsetfigurename{GraphTheory_Applied_}

В предыдущих разделах мы определили полугруппы, моноиды, группы, полукольца и кольца как <<честные>> алгебраические структуры: сформулировали аксиомы, которым они должны удовлетворять, и привели примеры.
На практике же при реализации алгоритмов анализа графов эти аксиомы часто соблюдаются лишь частично, а на первый план выходят вопросы представления данных, эффективности операций и удобства программного интерфейса.
Данный раздел посвящён именно таким прикладным аспектам~\sidecite{7761646}.

Начнём с простого наблюдения.
С точки зрения программиста естественно уметь выполнять матричное умножение в ситуации, когда аргументы и результат имеют разный тип.
Пусть даны матрицы
\[
    A: \texttt{Matrix}\langle T_1 \rangle, \qquad
    B: \texttt{Matrix}\langle T_2 \rangle,
\]
а результирующая матрица должна иметь тип
\[
    C: \texttt{Matrix}\langle T_3 \rangle.
\]
Для реализации такого обобщённого умножения потребуется определить всего три сущности:
\[
    \otimes: T_1 \times T_2 \to T_3, \qquad
    \oplus: T_3 \times T_3 \to T_3, \qquad
    \Bbbzero \in T_3,
\]
где $\otimes$~--- операция <<умножения>> ячеек, $\oplus$~--- операция <<сложения>> (агрегации) частичных результатов, а $\Bbbzero$~--- нейтральный элемент $\oplus$, используемый, например, для инициализации результирующей матрицы.
Само определение умножения при этом остаётся неизменным:
\[
    C[i,j] = \bigoplus_{k=0}^{n-1} A[i,k] \otimes B[k,j].
\]
Никаких дополнительных требований (ассоциативность, коммутативность, дистрибутивность) к операциям $\otimes$ и $\oplus$ не предъявляется~--- достаточно, чтобы их сигнатуры были согласованы с типами аргументов и результата%
\sidenote{Речь идёт не просто о перегрузке операторов, а о полноценных пользовательских функциях, передаваемых в качестве аргументов обобщённой процедуре умножения.}.
Так, в сообществе GraphBLAS термином <<полукольцо>> могут называть конструкцию, у которой типы аргументов умножения различны, что формально полукольцом в алгебраическом смысле не является; такая вольность оправдана практическими потребностями и устоялась в литературе%
\sidenote{
    За деталями обсуждения терминологии мы отсылаем читателя к~\cite{7761646}, а также~\cite{Garbar_Erin_Chernikov_Panfilyonok_Grigorev_2026}.
}.

При определении матрицы смежности (раздел~\ref{sec:GraphTheoryBasicDefinitions}) мы ввели множество $\Opt{L} = \{\Bbbzero\} \cup L$, где $\Bbbzero$ означает <<ребро отсутствует>>.
В языках программирования такой конструкции естественно соответствует тип \texttt{Option} (OCaml, F\#) или \texttt{Maybe} (Haskell): значение \texttt{Some x} означает <<есть ребро с меткой \texttt{x}>>, а \texttt{None}~--- <<ребра нет>>.
Таким образом, матрица смежности графа с метками типа \texttt{Lbl} имеет тип $\texttt{Matrix}\langle\texttt{Option}\langle\texttt{Lbl}\rangle\rangle$
\sidenote{%
    Можно показать, что в данном контексте $\texttt{Option}$ естественным образом образует моноид: $\texttt{None}$ играет роль $\Bbbzero$, а композиция двух значений \texttt{Some} даёт \texttt{Some} (объединение, сложение или иная операция, определённая над $L$).
}.
В реальных графах число рёбер, как правило, существенно меньше $|V|^2$, поэтому матрицы смежности разрежены.
Способы компактного хранения таких матриц подробно рассмотрены в разделе~\ref{sec:SparseMatricesAndVectors}.
Здесь же отметим ключевой для наших целей факт: разреженная матрица хранит только <<непустые>> ячейки (те, что соответствуют \texttt{Some x}), а ячейки со значением \texttt{None} восстанавливаются неявно и не хранятся.
Это порождает известную в сообществе GraphBLAS \emph{проблему явных и неявных нулей}~\sidecite{7761646}: комбинация двух не-хранимых ячеек не должна порождать новую хранимую ячейку.

Перейдём к операциям над разреженными матрицами.
\emph{Поэлементная операция} над двумя матрицами~--- это операция, которая применяется независимо к каждой паре ячеек с одинаковыми индексами.
Типизация такой операции имеет вид
\[
    \texttt{Option}\langle T_1\rangle \to \texttt{Option}\langle T_2\rangle \to \texttt{Option}\langle T_3\rangle.
\]
Заметим, что эта сигнатура допускает определение операции, возвращающей \texttt{Some x} даже тогда, когда оба аргумента равны \texttt{None}, что при работе с разреженными матрицами нежелательно%
\sidenote{Данная проблема подробно разбирается в~\cite{Garbar_Erin_Chernikov_Panfilyonok_Grigorev_2026}, где предложено решение на уровне системы типов.}.
Тем не менее, при дисциплинированном использовании эта сигнатура вполне работоспособна, и именно на её основе можно построить универсальную поэлементную операцию $\texttt{map2}$~\sidecite{7761646}:
\[
    \texttt{map2}\ (\texttt{op}: \texttt{Option}\langle T_1\rangle \to \texttt{Option}\langle T_2\rangle \to \texttt{Option}\langle T_3\rangle)\ (\texttt{m}_1, \texttt{m}_2) = \texttt{result}.
\]
Одна и та же функция $\texttt{map2}$, параметризованная разными операциями $\texttt{op}$, реализует и поэлементное сложение, и поэлементное умножение, и маскирование.

Рассмотрим конкретные примеры для целых чисел.

Операция поэлементного сложения:
\begin{align*}
    \texttt{op\_int\_add}\ (\texttt{Some}\ x)\ (\texttt{Some}\ y) &= \texttt{Some}\ (x + y), \\
    \texttt{op\_int\_add}\ (\texttt{Some}\ x)\ \texttt{None} &= \texttt{Some}\ x, \\
    \texttt{op\_int\_add}\ \texttt{None}\ (\texttt{Some}\ y) &= \texttt{Some}\ y, \\
    \texttt{op\_int\_add}\ \texttt{None}\ \texttt{None} &= \texttt{None}.
\end{align*}
Если в результате сложения получился <<честный>> ноль, можно дополнительно вернуть \texttt{None}, чтобы не сохранять его.

Операция поэлементного умножения, напротив, возвращает непустое значение только тогда, когда оба аргумента непусты:
\begin{align*}
    \texttt{op\_int\_mult}\ (\texttt{Some}\ x)\ (\texttt{Some}\ y) &= \texttt{Some}\ (x * y), \\
    \texttt{op\_int\_mult}\ \_\ \_ &= \texttt{None}.
\end{align*}

Маскирование~--- ещё один важный частный случай $\texttt{map2}$.
\emph{Структурная маска} оставляет значения первой матрицы только там, где во второй матрице есть непустые значения:
\[
    \texttt{op\_mask}\ (\texttt{Some}\ x)\ (\texttt{Some}\ y) = \texttt{Some}\ x,
\]
и \texttt{None} в остальных случаях.

\emph{Инвертированная структурная маска} действует наоборот: оставляет значения первой матрицы только там, где во второй матрице ячейка пуста.
Такая операция используется, например, в алгоритме BFS (раздел~\ref{sec:BFS}) при формировании нового фронта с учётом уже посещённых вершин:
\[
    \texttt{op\_inv\_mask}\ (\texttt{Some}\ x)\ \texttt{None} = \texttt{Some}\ x,
\]
и \texttt{None} в остальных случаях.

Наконец, можно определить \emph{предикатную маску}: оставить значение первой матрицы, если значение второй матрицы (когда оно есть) удовлетворяет заданному предикату $p$.
Структурная и инвертированная маски являются её частными случаями~--- с предикатами <<аргумент непуст>> и <<аргумент пуст>> соответственно.

Идея параметризации операций естественным образом переносится и на матричное умножение.
В терминах типов этому соответствует функция $\texttt{mxm}$:
\begin{align*}
    \texttt{mxm}\ & : \\
    & \quad \texttt{m}_1: \texttt{Matrix}\langle\texttt{Option}\langle T_1\rangle\rangle \\
    & \to \texttt{m}_2: \texttt{Matrix}\langle\texttt{Option}\langle T_2\rangle\rangle \\
    & \to \texttt{op\_mult}: \texttt{Option}\langle T_1\rangle \to \texttt{Option}\langle T_2\rangle \to \texttt{Option}\langle T_3\rangle \\
    & \to \texttt{op\_add}: \texttt{Option}\langle T_3\rangle \to \texttt{Option}\langle T_3\rangle \to \texttt{Option}\langle T_3\rangle \\
    & \to \texttt{zero}: \texttt{Option}\langle T_3\rangle \\
    & \to \texttt{result}: \texttt{Matrix}\langle\texttt{Option}\langle T_3\rangle\rangle.
\end{align*}
Здесь $\texttt{op\_mult}$ задаёт, как перемножить ячейку первой матрицы с ячейкой второй, $\texttt{op\_add}$~--- как сложить несколько частичных произведений для одной позиции результата, а $\texttt{zero}$~--- нейтральный элемент $\texttt{op\_add}$.

Вычисление происходит по формуле
\[
    \texttt{result}[i,j] = \texttt{op\_add}_{k}\ \texttt{op\_mult}\ \texttt{m}_1[i,k]\ \texttt{m}_2[k,j],
\]
где $\texttt{op\_add}_k$ обозначает итеративное применение $\texttt{op\_add}$ по $k$, начиная с $\texttt{zero}$.

Отметим важное различие с классическим определением матричного умножения над полукольцом (раздел~\ref{sec:PathAlgebra}).
В рамках описанного подхода $\texttt{op\_mult}$ и $\texttt{op\_add}$ не обязаны удовлетворять аксиомам ассоциативности, коммутативности или дистрибутивности.
На практике при решении задач анализа графов используемые операции часто не образуют полукольцо в математическом смысле, но это не мешает им быть полезными: достаточно, чтобы они были согласованы с конкретной решаемой задачей.
Когда же операции действительно образуют полукольцо, $\texttt{mxm}$ в точности совпадает с умножением матриц над этим полукольцом.

При решении некоторых задач анализа графов поэлементным операциям бывает полезно иметь доступ к координатам обрабатываемых ячеек~--- фактически, к номерам вершин.
В этом случае сигнатура операции <<умножения>> расширяется:
\[
    \otimes: (\texttt{int} \times \texttt{int} \times T_1) \times (\texttt{int} \times \texttt{int} \times T_2) \to T_3,
\]
где первые два аргумента каждого кортежа~--- номера строки и столбца обрабатываемой ячейки.
Эта возможность особенно полезна, когда результат операции зависит от того, между какими именно вершинами находится ребро.

Рассмотрим ещё один полезный на практике приём~--- выделение подмножества рёбер с использованием диагональных матриц.

Умножение матрицы смежности $M$ на диагональную матрицу $D$ справа позволяет выбрать только те рёбра, конечная вершина которых принадлежит заданному подмножеству $S \subseteq V$ (диагональ $D[i,i] = 1 \iff i \in S$).
Иными словами, $M \times D$ оставляет входящие в $S$ рёбра.

\begin{example}
    Рассмотрим граф $\mbfscrG_1$ из примера~\ref{exmpl:graph_visualization}.
    Выделим входящие рёбра для вершин $S = \{0, 2\}$ (на рисунке~\ref{fig:incomingEdges} выделены синим и красным цветами соответственно).
    Умножение справа на диагональную матрицу $D = \operatorname{diag}(1, 0, 1, 0)$ даёт:

    \begin{marginfigure}
    \begin{center}
        \begin{tikzpicture}[shorten >=1pt,auto]
  \node[state, fill=blue!20] (q_0)                      {$0$};
  \node[state] (q_1) [above right=of q_0]               {$1$};
  \node[state, fill=red!20] (q_2) [right=of q_0]        {$2$};
  \node[state] (q_3) [right=of q_2]                     {$3$};
  \path[->]
    (q_0) edge  node {} (q_1)
    (q_1) edge  node {} (q_2)
    (q_2) edge  node {} (q_0)
    (q_2) edge[bend left, above]  node {} (q_3)
    (q_3) edge[bend left, below]  node {} (q_2);
  {
    \path[->, red, thick]
    (q_1) edge node {} (q_2)
    (q_2) edge  node {} (q_0)
    (q_3) edge[bend left, below]  node {} (q_2);
  }
\end{tikzpicture}

    \end{center}
    \caption{Фильтрация входящих рёбер для вершин 0 и 2}
    \label{fig:incomingEdges}
\end{marginfigure}

    \[
        \left(
        \begin{array}{>{\columncolor{blue!20}}cc>{\columncolor{red!20}}cc}
            0 & 1 & 0 & 0 \\
            0 & 0 & 1 & 0 \\
            1 & 0 & 0 & 1 \\
            0 & 0 & 1 & 0
        \end{array}\right)
        \times
        \left(
        \begin{array}{>{\columncolor{blue!20}}cc>{\columncolor{red!20}}cc}
            1 & 0 & 0 & 0 \\
            0 & 0 & 0 & 0 \\
            0 & 0 & 1 & 0 \\
            0 & 0 & 0 & 0
        \end{array}\right)
        =
        \left(
        \begin{array}{cccc}
            0 & 0 & 0 & 0 \\
            0 & 0 & \cellcolor{red!20}{1} & 0 \\
            \cellcolor{blue!20}{1} & 0 & 0 & 0 \\
            0 & 0 & \cellcolor{red!20}{1} & 0
        \end{array}\right).
    \]
    Выделенные цветом ячейки соответствуют рёбрам: $(2,0)$~--- входящее в вершину $0$, и $(1,2)$, $(3,2)$~--- входящие в вершину $2$.
\end{example}

Умножение слева на диагональную матрицу $D \times M$, напротив, выбирает рёбра, исходящие из вершин заданного подмножества $S$:
\[
    (D \times M)[i,j] = D[i,i] \times M[i,j].
\]

\begin{example}
    Для того же графа и того же подмножества $S = \{0, 2\}$ умножение слева даёт фильтрацию исходящих рёбер (рисунок~\ref{fig:outgoingEdges}):

\begin{marginfigure}
    \begin{center}
        \input{part_01_Prep/chapter_03_GraphTheoryIntro/figures/04_AppliedAspects/outgoing_edges.tex}
    \end{center}
    \caption{Фильтрация исходящих рёбер для вершин 0 и 2 }
    \label{fig:outgoingEdges}
\end{marginfigure}


    \[
        \left(
        \begin{array}{cccc}
            \rowcolor{blue!20}
            1 & 0 & 0 & 0 \\
            0 & 0 & 0 & 0 \\
            \rowcolor{red!20}
            0 & 0 & 1 & 0 \\
            0 & 0 & 0 & 0
        \end{array}\right)
        \times
        \left(
        \begin{array}{cccc}
            \rowcolor{blue!20}
            0 & 1 & 0 & 0 \\
            0 & 0 & 1 & 0 \\
            \rowcolor{red!20}
            1 & 0 & 0 & 1 \\
            0 & 0 & 1 & 0
        \end{array}\right)
        =
        \left(
        \begin{array}{cccc}
            0 & \cellcolor{blue!20}{1} & 0 & 0 \\
            0 & 0 & 0 & 0 \\
            \cellcolor{red!20}{1} & 0 & 0 & \cellcolor{red!20}{1} \\
            0 & 0 & 0 & 0
        \end{array}\right).
    \]
    Выделены ячейки, соответствующие рёбрам: $(0,1)$~--- исходящее из вершины $0$, и $(2,0)$, $(2,3)$~--- исходящие из вершины $2$.
\end{example}

Разумеется, описанные выше типобезопасные абстракции ($\texttt{Option}$, $\texttt{map2}$, $\texttt{mxm}$) являются идеализацией.
В реальном стандарте GraphBLAS~\sidecite{7761646}~--- открытом стандарте, описывающем набор примитивов и операций разреженной линейной алгебры для построения алгоритмов анализа графов,~--- используется более <<плоское>> представление данных (без типа $\texttt{Option}$ в сигнатурах), однако ключевые идеи сохраняются: операции параметризуются пользовательскими полукольцами и моноидами, а поэлементным операциям доступны индексы обрабатываемых ячеек%
\sidenote{
    С точки зрения строгой математической терминологии именование произвольной параметризованной операции <<полукольцом>> или <<моноидом>> не вполне корректно.
    Однако идея близка, и для практических целей сообщество GraphBLAS не изобретает новых терминов, пользуясь устоявшимися названиями с некоторой вольностью.
    Детальное обсуждение этого вопроса можно найти в~\cite{7761646,Garbar_Erin_Chernikov_Panfilyonok_Grigorev_2026}.
}.
Стандарт позволяет единообразно выражать широкий класс графовых алгоритмов~--- от поиска кратчайших путей (тропическое полукольцо) до вычисления достижимости (булево полукольцо).

На практике этот стандарт реализуется в нескольких библиотеках: наиболее известная из них~--- SuiteSparse:GraphBLAS~\sidecite{Davis2018Algorithm9S}, широко используемая в высокопроизводительных графовых вычислениях.
Активно развиваются и GPGPU-реализации, такие как GraphBLAST.

Независимо от конкретной реализации, главное преимущество GraphBLAS-подхода для программиста~--- возможность выражать графовые алгоритмы через композицию матричных операций, делегируя заботу о разреженности, параллелизме и низкоуровневых оптимизациях библиотеке.
Операции $\texttt{map2}$ и $\texttt{mxm}$, параметризованные пользовательскими функциями, являются естественным обобщением этого подхода: они позволяют отказаться от жёсткого требования задавать полукольцо и вместо этого указывать ровно те операции, которые нужны для конкретной задачи, с явными гарантиями корректности на уровне типов~\sidecite{7761646}.
